\scanpage{49}
Rabiots pour l'exposé précédent:

a) \deleted{\unsure{Définition}}\unsure{Déscription|Description} de la filtration canonique d'un F-V-cristal dans
le cas d'une base \inserted{un anneau \uline{parfait}} \handwritten{(}de car. $p$\handwritten{)}: elle est définie par un sous-module
$\mathrm{Fil}^{1}\ M_{\unsure{o|0}}$ \deleted{($M$ considéré comme module localement libre de type fini sur
$W=\underline W(A)$)} tel que l'on ait
\[
\mathrm{Fil}^{1\ (p)}=\mathrm{Ker}\ F_M=\mathrm{Im}\ V_M\quad ,
\]
ce qui le définit de façon unique grâce à l'hypothèse sur $A$ (NB tout
\deleted{\unsure{c|M}}Cristal de Dieudonné est alors "admissible").

b) Iso\deleted{cristaux}groupes \inserted{de BT} (devrait faire un sous-paragraphe). Définition
sur une base quelconque: on localise \deleted{pour $Z$} en remplaçant
$\mathrm{Hom}\ (G,H)$ par $\mathrm{Hom}(G,H)\otimes_{\underline Z_p}\underline Q_p=\mathrm{Hom}\otimes_{\underline Z}\underline Q$. (Définition provisoire; elle
\deleted{\uline{Proposition}} a l'inconvénient de ne pas fournir un champ, et devrait
peut-être être reconsidérée ...). On définit de même les isocristaux
de Dieudonn\deleted{\unsure{è}}é, et on obtient donc un foncteur
\[
\mathrm{IsoBT}(S)\qquad\qquad\mathrm{Iso}\deleted{\mathrm{cris}}\text{F-V-cris}(S)
\]
($p$ localement nilpotent sur $\deleted{\unsure{J}}S$). On trouve

\uline{Proposition} (prouvée plus tard): \deleted{Supposons $S$ quasi-compact, et soit
soit \unsure{$S'$|$S$|$S_0$} un \unsure{nilidéal d'id|idéal d'id}} \inserted{Soit $S\qquad S'$ une immersion nilpoten}te d'Idéal $J$ annulé par une puissance de
$J$, alors le foncteur restriction
\[
\mathrm{IsoBT}(S')\handwritten{\longrightarrow}\mathrm{IsoBT}(S)
\]
est une équivalence de catégories.

On en conclut l'énoncé VII 10.8.2 de mes notes de Pise.\inserted{\handwritten{Expliciter le cas d'un anneau de val. discrète d'inégales car.}}
\margin{\handwritten{\uline{\unsure{Th.|Th}}}}
On peut se poser le \deleted{question de l'image essentielle}

\uline{Problème} 4. Quel est l'image essentielle dudit foncteur ? Suppposant
$A$ noethérien séparé et complet pour la topologique $p$-adique pour fixer
les idées, y a-t-il un critère valuatif d'appartenance à l'image essentielle ? Si $A$ est un anneau de val. discrète \deleted{\unsure{(}}d'inégales caractéristiques, l'image essentielle est elle un ouvert de \deleted{\unsure{z}}Zariski de la grassma-
% archive sheet number: 48

